% !TEX root = ../main.tex
\chapter{Op Emitters}
\label{ch:op_emitters}

The op emitters are the per-operation translation functions that
convert a TraceIR node into one or more \code{MPSGraphTensor}
operations. The MpsBuilder (\chref{ch:mps_builder}) dispatches
to them by op name; each emitter is responsible for one op's
translation. \Lucid\ 3.5 ships 183 emitters covering the common
operations of training workloads.

This chapter describes the emitter signature, the per-category
emitter patterns, the rules for adding a new emitter, and the
small subset of ops whose emitters are non-trivial.

\section{The Emitter Signature}
\label{sec:op_emitters:signature}

Each emitter is a function with the canonical signature:

\begin{lstlisting}[style=lucidcpp,language=C++]
std::vector<MPSGraphTensor*> emit_<op>(
    MPSGraph* graph,
    const std::vector<MPSGraphTensor*>& inputs,
    const std::vector<Attribute>& attrs);
\end{lstlisting}

The emitter receives the \MPSGraph\ being built, the input
tensor placeholders (already in the graph), and the op's
attributes (scalar parameters). It returns the output tensor
placeholders.

\subsection{Pure translation}
\label{ssec:op_emitters:pure_translation}

An emitter does not execute any actual computation. It only
inserts \code{MPSGraphTensor} nodes into the graph; \MPSGraph's
compiler later turns the graph into kernels.

This makes emitters simple: a few lines per op in the typical
case, calling one \MPSGraph\ method with the appropriate
arguments.

\section{Categories of Emitters}
\label{sec:op_emitters:categories}

The 183 emitters group naturally.

\subsection{Elementwise}
\label{ssec:op_emitters:elementwise}

The 60+ elementwise ops have nearly identical emitter
structure:

\begin{lstlisting}[style=lucidcpp,language=C++]
std::vector<MPSGraphTensor*> emit_add(
    MPSGraph* graph,
    const std::vector<MPSGraphTensor*>& inputs,
    const std::vector<Attribute>& attrs) {
  MPSGraphTensor* out = [graph additionWithPrimaryTensor:inputs[0]
                                         secondaryTensor:inputs[1]
                                                    name:nil];
  return {out};
}
\end{lstlisting}

The pattern: call the \MPSGraph\ method, return the output.
\Lucid\ generates these from a macro in \code{OpEmitters/Elementwise.mm}
to avoid the boilerplate.

\subsection{Reductions}
\label{ssec:op_emitters:reductions}

Reductions take an axis attribute:

\begin{lstlisting}[style=lucidcpp,language=C++]
std::vector<MPSGraphTensor*> emit_sum(
    MPSGraph* graph,
    const std::vector<MPSGraphTensor*>& inputs,
    const std::vector<Attribute>& attrs) {
  AxisSet axes = attrs[0].as_axis_set();
  bool keepdims = attrs[1].as_bool();
  MPSGraphTensor* out = [graph reductionSumWithTensor:inputs[0]
                                                 axes:to_nsarray(axes)
                                                 name:nil];
  if (!keepdims) {
    out = [graph squeezeTensor:out axes:to_nsarray(axes) name:nil];
  }
  return {out};
}
\end{lstlisting}

The \code{keepdims=False} case requires an explicit squeeze
because \MPSGraph's reduction always keeps the reduced dimensions
as size 1.

\subsection{Matrix multiplication}
\label{ssec:op_emitters:matmul}

\begin{lstlisting}[style=lucidcpp,language=C++]
std::vector<MPSGraphTensor*> emit_matmul(
    MPSGraph* graph,
    const std::vector<MPSGraphTensor*>& inputs,
    const std::vector<Attribute>& attrs) {
  MPSGraphTensor* out = [graph
      matrixMultiplicationWithPrimaryTensor:inputs[0]
                          secondaryTensor:inputs[1]
                                     name:nil];
  return {out};
}
\end{lstlisting}

\MPSGraph's matmul handles broadcasting and the batched form;
the emitter is one call.

\subsection{Convolution}
\label{ssec:op_emitters:convolution}

The convolution emitter is the most complex of the regular cases
because of the descriptor object that \MPSGraph\ uses:

\begin{lstlisting}[style=lucidcpp,language=C++]
std::vector<MPSGraphTensor*> emit_conv2d(
    MPSGraph* graph,
    const std::vector<MPSGraphTensor*>& inputs,
    const std::vector<Attribute>& attrs) {
  auto desc = [MPSGraphConvolution2DOpDescriptor descriptor];
  desc.strideInX = attrs[0].as_int();
  desc.strideInY = attrs[1].as_int();
  desc.paddingLeft = attrs[2].as_int();
  desc.paddingRight = attrs[3].as_int();
  desc.paddingTop = attrs[4].as_int();
  desc.paddingBottom = attrs[5].as_int();
  desc.dataLayout = MPSGraphTensorNamedDataLayoutNCHW;
  desc.weightsLayout = MPSGraphTensorNamedDataLayoutOIHW;

  MPSGraphTensor* out = [graph
      convolution2DWithSourceTensor:inputs[0]
                       weightsTensor:inputs[1]
                          descriptor:desc
                                name:nil];
  if (inputs.size() == 3) {
    // Bias provided
    out = [graph additionWithPrimaryTensor:out
                          secondaryTensor:inputs[2]
                                      name:nil];
  }
  return {out};
}
\end{lstlisting}

The descriptor pattern is shared across \MPSGraph's conv,
pooling, and BatchNorm emitters.

\subsection{LayerNorm}
\label{ssec:op_emitters:layernorm}

LayerNorm has a dedicated \MPSGraph\ primitive that the emitter
calls:

\begin{lstlisting}[style=lucidcpp,language=C++]
std::vector<MPSGraphTensor*> emit_layer_norm(
    MPSGraph* graph,
    const std::vector<MPSGraphTensor*>& inputs,
    const std::vector<Attribute>& attrs) {
  MPSGraphTensor* mean = [graph meanOfTensor:inputs[0]
                                       axes:[axes]
                                       name:nil];
  MPSGraphTensor* variance = [graph varianceOfTensor:inputs[0]
                                           meanTensor:mean
                                                 axes:[axes]
                                                 name:nil];
  MPSGraphTensor* normalized = [graph normalizationWithTensor:inputs[0]
                                                   meanTensor:mean
                                               varianceTensor:variance
                                                  gammaTensor:inputs[1]
                                                   betaTensor:inputs[2]
                                                      epsilon:attrs[0].as_float()
                                                         name:nil];
  return {normalized};
}
\end{lstlisting}

The fused \code{normalizationWithTensor} avoids materialising
the intermediate per-element normalisation, which is what makes
the \MPSGraph\ LayerNorm faster than the equivalent MLX
composite (\chref{ch:mpsgraph_dispatch}).

\subsection{Multi-head attention}
\label{ssec:op_emitters:mha}

The attention emitter calls \MPSGraph's
\code{scaledDotProductAttentionWithQueryTensor:keyTensor:valueTensor:scale:}
primitive when the mask is supported (causal or none); falls back
to the composite path otherwise.

This is the most performance-critical emitter: a Transformer
training step spends a third of its time inside attention. The
fused kernel saves both compute and memory.

\section{Adding a New Emitter}
\label{sec:op_emitters:adding}

The protocol for adding a new emitter:

\begin{enumerate}
  \item Locate the appropriate file in
    \code{lucid/\_C/compile/OpEmitters/} (by op category).
  \item Write the emitter function following the canonical
    signature.
  \item Add the self-registering struct at the bottom of the
    file.
  \item Add a parity test that compiles a single-op model
    containing the new op, runs both compiled and eager, and
    verifies they match.
  \item Rebuild; the new emitter is now in the table.
\end{enumerate}

Typical implementation cost: a few dozen lines for a simple
emitter, a few hundred for complex ones (attention,
normalisation families). The parity test is the gating
requirement; without it, the new emitter is silently broken.

\section{Coverage and Gaps}
\label{sec:op_emitters:coverage}

\tabref{tab:op_emitters:coverage} shows the per-category
coverage as of 3.5.

\begin{table}[t]
\centering
\small
\begin{tabular}{lrr}
\toprule
Category & Ops in registry & Emitters \\
\midrule
Elementwise unary    & 64 & 64 \\
Elementwise binary   & 52 & 52 \\
Reductions           & 22 & 22 \\
Shape / view         & 18 & 18 \\
Indexing             & 12 & 9  \\
Convolution / pool   & 14 & 14 \\
Normalisation        & 6  & 6  \\
Attention            & 3  & 2  \\
Linalg               & 31 & 0  \\
FFT                  & 22 & 0  \\
NN composites        & 12 & 8  \\
Einops               & 6  & 6  \\
Misc                 & 36 & 32 \\
\midrule
Total                & 298 & 233 (78\%) \\
\bottomrule
\end{tabular}
\caption[Emitter coverage.]{Per-category coverage of the op
emitter table. Linalg (31 ops) and FFT (22 ops) have no emitters
because \MPSGraph\ does not expose the underlying primitives;
those ops always fall back to the eager linalg-carveout path
(\chref{ch:backend_dispatch}). Coverage will grow with future
\MPSGraph\ releases.}
\label{tab:op_emitters:coverage}
\end{table}

\subsection{Why linalg has no emitters}
\label{ssec:op_emitters:linalg_gap}

\MPSGraph\ does not provide LAPACK-equivalent operations. The
linalg ops (\code{cholesky}, \code{qr}, \code{svd}, etc.)
therefore have no emitters, and the compile path falls back to
eager for any subgraph containing one. This matches the
linalg carve-out at the dispatcher level
(\chref{ch:backend_dispatch}).

\subsection{Why FFT has no emitters}
\label{ssec:op_emitters:fft_gap}

\MPSGraph\ added FFT primitives in macOS~14, but the integration
with autograd in the compile path is non-trivial (FFT's backward
involves Hermitian symmetry that the current emitter
infrastructure does not handle elegantly). Deferred.

\section{Generation vs Hand-Writing}
\label{sec:op_emitters:gen_vs_hand}

The elementwise emitters (60+ unary, 50+ binary) are nearly
identical and could be code-generated from the op registry.
\Lucid\ does generate them, via a build-time script
\code{tools/gen\_emitters.py} that reads the registry and emits
\code{OpEmitters/Elementwise.gen.mm}.

Non-elementwise emitters are hand-written. The 50-or-so
hand-written emitters each take a few hours to write, test, and
benchmark; over the course of the compile path's development,
this was the dominant implementation cost.

\section{Numerical Equivalence}
\label{sec:op_emitters:numerical_equiv}

Each emitter is parity-tested against its eager equivalent: the
parity test runs the same op on the same input through both
paths and checks that outputs agree to within FP32 tolerance
($10^{-5}$ relative, $10^{-4}$ absolute typically).

The fused \MPSGraph\ implementations sometimes differ from the
eager composites by a few ULPs (different reduction order, or
fused intermediate computation in FP32 vs the eager FP16). These
differences are documented per-op and the parity tolerance
adjusted accordingly.

\section{Performance of Compiled vs Eager}
\label{sec:op_emitters:perf}

Per-op timings on M3~Max for the most common ops, comparing
eager dispatch to compiled-graph dispatch:

\begin{table}[t]
\centering
\small
\begin{tabular}{lrrr}
\toprule
Op & Shape & Eager (ms) & Compiled (ms) \\
\midrule
matmul    & $(B, T, D) \times (D, D), B = 8, T = 1024, D = 768$ & 1.6 & 1.4 \\
gelu      & $(B, T, D) = (8, 1024, 3072)$                       & 0.7 & 0.5 \\
layernorm & $(B, T, D) = (8, 1024, 768)$                        & 0.8 & 0.2 \\
softmax   & $(B, H, T, T) = (8, 12, 1024, 1024)$                & 3.4 & 2.8 \\
sdpa      & MHA over same shape                                  & 11  & 2.6 \\
embedding & $(B, T) = (8, 1024)$, vocab 50000, D = 768           & 0.4 & 0.3 \\
\bottomrule
\end{tabular}
\caption[Per-op compile speedup.]{Per-op timings on M3~Max for
the most performance-critical ops. The speedup is per-op
(matmul: 1.1$\times$); much of the compile-path benefit comes
from \emph{inter-op fusion} (LayerNorm avoids materialising
intermediates, SDPA fuses the entire attention into one kernel).}
\label{tab:op_emitters:perf}
\end{table}

The per-op speedup is modest; the inter-op fusion is where the
compile path's real benefit comes from.

\section{Summary and Forward References}
\label{sec:op_emitters:summary}

The op emitter table is the per-operation translation layer
that turns a TraceIR into an \MPSGraph. \Lucid\ 3.5 ships 183
emitters covering 78\% of the registered ops; the remaining 22\%
(notably linalg and FFT) fall back to eager. The elementwise
emitters are code-generated from a template; the non-trivial
emitters (attention, layer norm, conv) are hand-written.

The next chapter, \chref{ch:fwd_bwd_compile}, treats the
forward+backward unification --- the design that lets the compile
path capture the entire training step as a single graph.
\chref{ch:compiled_module} treats the user-facing API.
\chref{ch:compile_hardening} treats the edge cases.

\begin{lucidnote}
For the user: the emitter table is invisible at the API surface.
Op coverage matters only when the user encounters an unsupported
op, in which case \code{strict=True} surfaces the name and the
user can route around it (compose with \code{where}, switch to
an equivalent supported op).
\end{lucidnote}
